\chapter{Background}
\label{ch:background}

This chapter introduces the two certification schemes, the family of
libraries under study, and the sec-certs project, and then reviews related
work.

\section{Common Criteria}
\label{sec:cc}

Common Criteria, standardized as ISO/IEC~15408~\cite{iso15408}, is a
framework for evaluating the security of IT products. The product under
evaluation is called the Target of Evaluation (TOE). The vendor describes
the TOE and its security claims in a Security Target (ST), which may claim
conformance to a Protection Profile, a reusable set of requirements for a
class of products. An accredited laboratory evaluates the TOE against those
claims at a chosen Evaluation Assurance Level (EAL), and a national scheme
issues the certificate and a certification report.

Three properties of the scheme matter for this thesis:

\begin{itemize}
  \item \textbf{The certificate covers the claims, not the whole product.}
    A library may be part of the TOE, part of its environment, or merely
    mentioned, so a mention alone does not establish that the library was
    evaluated.
  \item \textbf{Disclosure is not uniform.} CC does not require a list of
    third-party components in a form suited to analysis. Vendors decide how
    much to reveal, and the public version of an ST may be sanitized.
  \item \textbf{Certificates change over time.} A certified product can be
    updated through assurance continuity, which produces maintenance
    reports. These documents are a potential source of evidence on library
    upgrades.
\end{itemize}

{Cite the CC portal or CCRA for mutual recognition and assurance
continuity. The portal could not be retrieved automatically while drafting,
so check these statements against it directly. Also check the current status of
ISO/IEC~15408-1:2022, because ISO's own pages were inconsistent about it, and
decide whether to cite the CC version used by most certificates in the dataset.}

\section{FIPS 140}
\label{sec:fips}

FIPS~140 is the US federal standard for the security requirements of
cryptographic modules~\cite{fips1402,fips1403}. Modules are tested by
accredited laboratories and validated under the Cryptographic Module
Validation Program (CMVP), which publishes a certificate and a Security
Policy for each module. In contrast to CC, the object of a FIPS validation
is the cryptographic module itself, so a library can appear either as the
validated module or as a component inside a larger one.

The standard is in the middle of a generational change that directly affects
lifecycle analysis. FIPS~140-3 was approved in March~2019 and took effect in
September of that year. CMVP began accepting FIPS~140-3 submissions in
September~2020 and stopped accepting FIPS~140-2 submissions for new
certificates in April~2022. In September~2026, the FIPS~140-2 validations were
scheduled to move to the Historical List, although CMVP continues to support
buying and using such modules for existing systems~\cite{nistTransition}.
{Confirm the move took place as scheduled and how sec-certs records it.
The NIST page gives both 21 and 22 September.}

\section{OpenSSL and its forks}
\label{sec:libraries}

OpenSSL is an open-source implementation of the TLS protocol and a
general-purpose cryptographic library. Its current release policy provides a
new release every six months, and from version~3.5, a long-term support
release every two years with at least five years of
support~\cite{opensslReleaseStrategy}. This cadence is the upstream reference
against which certification lag is measured in Chapter~\ref{ch:method}.

OpenSSL has its own relationship with FIPS~140. Historically, FIPS support was
delivered through the separate OpenSSL FIPS Object Module, while the 3.x series
provides it through the FIPS Provider. The FIPS Provider illustrates why the
version numbers require care. Certificate \#4282 was first validated in
August~2022 and now lists versions 3.0.8 and 3.0.9~\cite{cmvp4282}. Versions
3.0.8 and 3.0.9 are also covered by the second FIPS~140-2 certificate, \#4811,
and version 3.1.2 by the FIPS~140-3 certificate \#4985~\cite{opensslFipsCve}.
According to an OpenSSL developer, the provider is designed to work with all
3.0.x releases~\cite{opensslUsers4282}. The version of the validated provider
can therefore differ from the version of the library that loads it.
{Add a primary source for the FIPS Object Module history.}

Throughout this thesis, a \emph{fork} is a code base with a documented
derivation from OpenSSL, and \emph{clone} is used only as the broader term from
the thesis assignment. Independent TLS libraries that do not derive from OpenSSL
are not treated as forks. Table~\ref{tab:libraries} lists the libraries in scope.

\begin{table}[htbp]
  \centering
  \caption{Libraries in scope and their relationship to OpenSSL}
  \label{tab:libraries}
  \begin{tabular}{lp{0.62\textwidth}}
    \toprule
    Library & Relationship to OpenSSL \\
    \midrule
    OpenSSL & The upstream project, including the FIPS Object Module and the
      FIPS Provider. \\
    LibreSSL & Forked in 2014 with the goal of modernizing the code base and
      improving security; developed mainly in the OpenBSD source
      tree~\cite{libressl}. \\
    BoringSSL & Google's fork, created to replace a growing set of private
      patches. Not intended for general use, and without API or ABI stability
      guarantees~\cite{langley2014boringssl,boringssl}. \\
    AWS-LC & Maintained by the AWS Cryptography team; based on code from both
      BoringSSL and OpenSSL~\cite{awslc}. \\
    quictls & A further member of the OpenSSL family of
      forks~\cite{stenberg2025forks}. \\
    \bottomrule
  \end{tabular}
\end{table}

The list is deliberately open. Identifying further forks, as well as
vendor-branded wrappers, is part of the first project phase
(Section~\ref{sec:heuristics}).

\section{The sec-certs project}
\label{sec:seccerts}

The sec-certs project~\cite{janovsky2024seccerts} collects CC and FIPS~140
certification artifacts, processes them, and publishes the results at
sec-certs.org. Its Python package exposes each scheme as a dataset object that
can be downloaded in processed form or loaded from JSON. Its processing
pipeline extracts data such as keywords from the documents and computes various
heuristics over the certificates~\cite{seccertsDocs}. The thesis builds on this
processed data instead of collecting the documents again.

\section{Related work}
\label{sec:related}

The closest related work is the sec-certs study by Janovsky et
al.~\cite{janovsky2024seccerts}. Its large-scale automated analysis of Common
Criteria certificates maps vulnerabilities from the National Vulnerability
Database to certified products using unsupervised models, and identifies the
security requirements associated with products that have fewer and less severe
vulnerabilities. This thesis uses the same data source but changes the unit of
analysis: it follows one family of cryptographic libraries across both schemes,
rather than following vulnerabilities across CC products.

Durumeric et al.~\cite{durumeric2014heartbleed} measured the population
vulnerable to Heartbleed, its patching over time, and exploitation attempts.
Their work observes deployed systems on the Internet. This thesis asks a
complementary question about the certified product population, which cannot be
observed by scanning alone and for which documents are the main evidence.

Ladisa et al.~\cite{ladisa2023sok} provide a taxonomy of attacks on open-source
software supply chains. It offers general context on the risks of inherited
dependencies, while this thesis focuses on measuring one such dependency rather
than on attack vectors.

{Extend with a systematic literature search: prior analyses of FIPS~140 and
CMVP data, software composition analysis, SBOM adoption, and LLM-based
information extraction from technical documents. Check whether anyone has already
studied OpenSSL in certified products before claiming this as a gap.}
